\chapter{Diskussion und Implikationen}

\section{Theoretische und wissenschaftliche Einordnung}

\subsection{Erste Modulationsannahme M1 – Willensbahnung}
Aus theoretischer Perspektive lässt sich das Ergebnis zur ersten Modulationsannahme (Szenario $XM_1$) von \cite{Kuhl2000The} mit der gezielten Berücksichtigung der affektiven Ausgangslage und der daraus abgeleiteten Interventionen erklären. Die Ergebnisse haben an dieser Stelle einen positiven Effekt von 9,99 Prozentpunkten der Experimentalgruppe gegenüber der Kontrollgruppe gezeigt. Während $B_K$ primär auf eine allgemeine traumasensible Gesprächsführung zurückgreift, berücksichtigt $B_{PSI}$ zusätzlich die zuvor identifizierte Inkongruenz zwischen Intentionsgedächtnis und intuitiver Verhaltenssteuerung. 

Durch Interventionen wie Bestärkung, Pendeltechnik und Ressourcenaktivierung wird gezielt die Erzeugung positiver Affekte angestrebt. Die Ergebnisse deuten darauf hin, dass durch den Einsatz dieser Interventionen eine Verschiebung des untersuchten linguistischen Indikators in die erwartete Richtung erreicht werden kann. Diese Ergebnisse bestärken somit die einen förderlichen Einsatz der Pendeltechnik entsprechend \cite{kuhl2005neue} sowie dem Einsatz im menschlichen Coaching \citep{kuhl2014psi}. Die Interventionen lassen sich relativ klar in sprachliche Impulse übersetzen, was die Steuerung und Verwendung im Modell erleichtert. An dieser Stelle sollte jedoch kritisch beleuchtet werden, dass die Ergebnisse vorrangig eine Veränderung linguistischer Indikatoren und nicht die tatsächliche Aktivierung der psychischen Systeme beim Menschen belegen. Es kann keine Aussage über eine reale Verbesserung der Handlungsfähigkeit oder Selbstregulation getroffen werden.

Die gezielte Berücksichtigung der Inkongruenz zwischen IG und IVS bezieht sich auf die Annahmen, dass das Intentionsgedächtnis Ziele und Absichten aufrecht hält, während die Umsetzung durch eine unzureichende Aktivierung der IVS gehemmt sein kann \citep{Kuhl2000The}. Durch die Aktivierung des IVS im Szenario werden persönliche Kompetenzen gefördert \citep{kuhl2004systemdiagnostik}. Positive Affekte können die Überführung von Absichten in konkretes Handeln unterstützen \citep{kuhl2005neue}. Die Ergebnisse des ersten Szenarios stehen somit im Einklang mit den theoretischen Annahmen der ersten Modulationsannahme.

Insgesamt verdeutlichen die Ergebnisse das Potenzial, psychologische Zustandsanalysen und daraus abgeleitete Interventionslogiken in KI-basierte Coaching-Systeme zu übertragen. Im Vergleich zu einer allgemeinen traumasensiblen Gesprächsführung scheint die zusätzliche Berücksichtigung der psychologischen Ausgangslage im Szenario $XM_1$ einen messbaren Mehrwert zu bieten. Damit greift die Arbeit die Bedeutung psychologischer und beziehungsbezogener Gestaltungsfaktoren für die Mensch-Chatbot-Interaktion auf, wie sie unter anderem von \cite{mai2020ai} thematisiert wird.


\subsection{Zweite Modulationsannahme M2 – Selbstentwicklung durch Affektregulation}
Bezüglich der zweiten Modulationsannahme von \cite{Kuhl2000The}, die auf Selbstentwicklung durch Affektregulation abzielt, zeigen die Ergebnisse des Versuchs keinen signifikanten Effekt zwischen kontrollbedingung $B_K$ und Experimentalbedingung $B_PSI$. Im direkten Vergleich hat die Experimentalgruppe sogar um -0,54 Prozentpunkte schlechter abgeschnitten. 

Eine mögliche Erklärung ergibt sich aus der unterschiedlichen Komplexität der beiden untersuchten Modulationsprozesse. Während die Aktivierung des IVS durch relativ konkrete motivationale Impulse, beispielsweise Ermutigung, Zielorientierung und Ressourcenaktivierung, unterstützt werden kann, setzt die Selbstentwicklung eine komplexere Verarbeitung emotionaler Erfahrungen voraus \citep{kuhl2005neue}. Das OES richtet die Aufmerksamkeit auf Fehler, Gefahren und negative Einzelaspekte. Die Aktivierung des EG setzt dagegen einen Zugang zu persönlichen Erfahrungen, Bedürfnissen und Werten voraus. Somit geht es geht nicht nur um eine Veränderung der Stimmung, sondern um die Integration emotionaler Erfahrungen in das Selbstsystem.

An dieser Stelle sollten die inhärenten Grenzen einer rein textbasierten Dialogsteuerung kritisch betrachtet werden. Eine sprachlich erzeugte Beruhigung ist nicht zwangsläufig mit einer tatsächlichen, neurobiologisch emotionalen Selbstberuhigung gleichzusetzen. Eine rein linguistisch messbares Nutzer-Sentiment im Interaktionsverlauf spiegelt somit nicht zwingend die reale Herabregulierung des negativen Affekts im Sinne der zweiten Modulationsannahme wider \citep{Kuhl2000The}. Im Sinne des Embodiment-Ansatzes nach \cite{tschacher2012bedeutung} bleibt eine nachhaltige psychische Veränderung ohne viszerale Resonanz und die Aktivierung somatischer Marker oft wirkungslos. 
Auf Grundlage der Ergebnisse könnte man annehmen, dass die generierte Antwort zu wenig auf die Erfahrungsintegration ausgerichtet sein könnte. Reine kognitive Einsicht oder Trost reichen nach \cite{kuhl2005neue} meist nicht aus, um traumabedingte Blockaden im Objekterkennungssystem (OES) aufzulösen, da das betroffene Selbstsystem primär über emotionale und erfahrungsbasierte Netzwerke angesprochen wird \citep{storch2007ressourcenorientiertes}. Für die Aktivierung des Extensionsgedächtnisses (EG) stellt eine reale Beziehungserfahrung einen wichtigen Faktor dar \citep{kuhl2014psi}. In dem Versuch stellt dies in der Mensch-Maschine-Interaktion eine essenzielle Herausforderung dar \citep{mai2020ai}. 
Psychische Selbstkompetenzen und die Fähigkeit zur Affektregulation können meist nur durch korrigierende Sicherheitserfahrungen in Beziehungen entwickelt werden \citep{Levine2010,kuhl2014psi}. An dieser Stelle stößt die simulationstechnische Erzeugung von psychologischer Sicherheit durch dialogische KI an eine Grenze \citep{mai2020ai}.

Das Ergebnis bedeutet nicht, dass die zweite Modulationsannahme grundsätzlich nicht durch KI unterstützt werden kann. Es konnte kein zusätzlicher Vorteil gegenüber der Kontrollbedingung nachgewiesen werden. o	Eine regelbasierte PSI-Steuerung reicht demnach möglicherweise nicht für alle psychischen Prozesse gleichermaßen aus.



\subsection{Vergleichende Gesamtbewertung der beiden Modulationsannahmen}
Im direkten Vergleich der beiden Szenarien zeigt sich, dass die PSI-basierte Steuerungslogik insbesondere bei der ersten Modulationsannahme einen deutlichen zusätzlichen Steuerungseffekt erzielt. Für die zweite Modulationsannahme konnte hingegen kein entsprechender Vorteil gegenüber der Kontrollgruppe nachgewiesen werden.

Die Ergebnisse verdeutlichen damit, dass die Übertragung der PSI-Theorie in die Entscheidungslogik eines KI-basierten Chatbots grundsätzlich Potenzial besitzt, jedoch nicht für sämtliche psychischen Prozesse gleichermaßen erfolgreich umgesetzt werden konnte.

Besonders hervorzuheben ist, dass die Kontrollgruppe in beiden Szenarien ebenfalls positive Systemverschiebungen aufweist. Dies deutet darauf hin, dass bereits eine allgemeine traumasensible Gesprächsführung bestimmte psychologische Aktivierungsprozesse unterstützen kann \citep{mai2023chatbots}. Die zusätzliche Integration der PSI-basierten Steuerungslogik scheint insbesondere dann einen Mehrwert zu bieten, wenn die angestrebte Veränderung durch konkrete motivationale Impulse unterstützt werden kann.


\subsection{Beantwortzung Forschungsfrage}
Für die Beantwortung der Forschungsfrage ergibt sich daraus, dass eine gezielte Verknüpfung von Zustandsanalyse, affektiver Bewertung und interventionsbasierter Dialogsteuerung ein relevantes Designmerkmal für KI-basierten Coaching-Systeme darstellt. Die Ergebnisse sprechen dafür, dass die explizite Berücksichtigung psychologischer Zustände eine differenziertere und zielgerichtetere Interaktion ermöglicht als eine ausschließlich reaktive Gesprächsführung.

Gleichzeitig muss berücksichtigt werden, dass die Untersuchung auf kontrollierten Dialogsimulationen und einer linguistischen Operationalisierung psychischer Systeme basiert. Die gemessenen Systemverschiebungen stellen daher keine unmittelbaren Veränderungen tatsächlicher psychischer Zustände von Personen dar. Insbesondere kann aus der Aktivierung sprachlicher Indikatoren für das Extensionsgedächtnis oder die intuitive Verhaltenssteuerung nicht unmittelbar auf eine tatsächliche psychologische Veränderung geschlossen werden.

Zusammenfassend lässt sich festhalten, dass die Evaluation die Funktionsfähigkeit des entwickelten Frameworks teilweise bestätigt. Für die erste Modulationsannahme konnte ein statistisch signifikanter und starker zusätzlicher Steuerungseffekt nachgewiesen werden. Für die zweite Modulationsannahme konnte die angenommene Überlegenheit hingegen nicht bestätigt werden. Das entwickelte Modell stellt damit einen konzeptionellen Ansatz dar, dessen Potenzial insbesondere hinsichtlich der Willensbahnung durch die vorliegenden Simulationsergebnisse gestützt wird. Eine weiterführende Untersuchung ist erforderlich, um die Übertragbarkeit der Ergebnisse auf reale Coaching-Interaktionen und die tatsächliche Unterstützung selbstregulatorischer Prozesse zu überprüfen.


\section{Praktische Implikationen}

\section{Ethische Implikationen und Safety}

\section{Kritische Betrachtung}

